%======================================================================
%  Supplementary Methods -- pretraining (paper section 2)
%
%  Source: /home/canall/DNA_LM/plasmids_GPN/manuscript/pretraining/
%    02_methods_corpus_and_clustering.md
%    03_methods_epoch_construction.md              (the canonical variant;
%       the _factorized_v1 and _factorized_DRAFT siblings are stale)
%    04_methods_windowing_and_tokenization.md
%    05_methods_architecture_objective_and_optimization.md
%  snapshot 2026-08-28. These files mirror the artifact's Methods band and
%  self-declare as the canonical source of truth for it.
%
%  Superscript citations map as in 02_pretraining.tex.
%======================================================================

\subsection*{Pretraining corpus and sequence clustering}

\subsubsection*{Source dataset}

\model was pretrained on plasmid sequences from the Integrated Microbial
Genomes / Plasmid Repository (IMG/PR), \cite{camargo2024img} comprising
699,973 sequences distributed across 214,950 Plasmid Taxonomic Units
(PTUs) \cite{redondosalvo2020pathways}. Sequences were used together with
their IMG/PR metadata, including GOLD ecosystem annotations
\cite{mukherjee2023twenty}, host taxonomy, assembly completeness, topology
annotation, and per-sequence composition statistics.

\subsubsection*{Filtering}

Four criteria were applied in sequence:

\begin{center}
\small
\begin{tabular}{clc}
\toprule
Step & Criterion & Sequences remaining \\
\midrule
--- & IMG/PR source & 699,973 \\
1 & length $\geq$ 2,048 bp & 699,110 \\
2 & ambiguous-base fraction $\leq$ 50\% and DUST-masked fraction $\leq$ 50\% & 699,085 \\
3 & sequences annotated as concatemers excluded & 693,644 \\
\bottomrule
\end{tabular}
\end{center}

The length threshold matches the model's input window, so that every retained
sequence yields at least one full-length training window. The composition
filters remove 25 sequences dominated by ambiguous bases or low-complexity
repeats. Concatemers --- assemblies containing multiple tandem copies of the
same replicon --- were excluded using the IMG/PR \texttt{topology} annotation,
removing 5,441 sequences; they were not re-detected computationally at this
stage.

\subsubsection*{Sequence clustering}

Pairwise sequence relationships were computed with skani
\cite{shaw2023fast}, which estimates the average nucleotide identity (ANI) ---
the mean identity across the aligned regions of two sequences --- together
with the alignment fraction (AF) --- the proportion of a sequence spanned by
that alignment. skani was run in all-versus-all (triangle) mode over the full
sequence set, treating each sequence as an individual genome, under its
high-sensitivity (\texttt{--slow}) preset with an 85\% ANI screening
threshold (\texttt{-s 85}) and a 95\% minimum alignment fraction
(\texttt{--min-af 95}); the exact invocation is recorded in the provenance
notes.

From these values, two levels of within-PTU clustering were defined by
independent criteria applied to the same sequence set, not by a single
agglomerative procedure.

\textbf{Dereplication clusters} were defined at $\geq$99\% ANI with $\geq$95\%
\emph{bidirectional} alignment fraction, requiring that both sequences in a
pair cover at least 95\% of each other's length. Because the criterion is
symmetric, only pairs that are near-identical in both sequence and length are
collapsed: a short fragment is not absorbed into a much longer plasmid it
happens to match, and so survives as an independent training example. This
level corresponds to repeated observation of the same molecule --- the same
strain or isolate.

\textbf{Sub-PTUs} were defined at $\geq$95\% ANI with $\geq$95\%
\emph{unidirectional} alignment fraction, requiring only that at least one
sequence of a pair is $\geq$95\% covered by the alignment. The identity
threshold follows the conventional boundary used to delimit prokaryotic
species from genome-wide ANI, \cite{jain2018high} though we note that species
concepts for plasmids remain operational rather than formally defined, and we
use the threshold as a practical delimiter of plasmid-level structure. The
unidirectional coverage criterion is deliberate: it groups a complete plasmid
with partial assemblies of the same plasmid, which is the common situation
for metagenomically derived sequences. It follows that sub-PTUs may contain
sequences of substantially different lengths, and that membership reflects
sequence correspondence rather than complete functional equivalence, since a
partial assembly may lack accessory regions present in a complete one.

Across the full database this yields 505,821 dereplication clusters and
343,291 sub-PTUs. PTUs contain a mean of 1.60 sub-PTUs and a median of one,
with a maximum of 737; 80.8\% of PTUs resolve into a single sub-PTU. Host
assignments span 49 phyla in the full collection.

\subsubsection*{Hierarchy validation}

Because the two clustering levels are defined independently, strict
containment is not guaranteed a priori, and we therefore measured it.
Containment is exact at the PTU level: no sub-PTU and no dereplication
cluster spans more than one PTU. This is the level at which the training and
held-out split is defined, so the split is unambiguous.

\subsubsection*{Training and held-out split}

The corpus was partitioned by host species: all PTUs containing a plasmid
whose host species was \emph{Shewanella oneidensis} or
\emph{Gluconacetobacter tumulisoli} were reserved as a species-disjoint
held-out set (3 PTUs, 6 sequences; 5 sub-PTU representatives, the longest
sequence per cluster). Splitting at the level of whole PTUs, rather than of
individual sequences, prevents leakage between the two partitions: because no
sub-PTU or dereplication cluster crosses a PTU boundary, no close relative of
a held-out sequence remains in training.

The training corpus is therefore 693,638 sequences spanning 213,393 PTUs,
340,697 sub-PTUs, and 500,466 dereplication clusters. Its largest PTU
contains 735 sub-PTUs, and its host assignments span 48 phyla, compared with
49 in the full collection. Sequence lengths have a mean of 24.0 kb and a
median of 10.1 kb.

\subsubsection*{PTU categories}

Each PTU in the training corpus was classified by which levels of the
hierarchy are populated within it:

\begin{itemize}
\item \textbf{Major} --- more than one sub-PTU (multiple plasmid sequences captured within the clade);
\item \textbf{Moderate} --- exactly one sub-PTU but more than one dereplication cluster (a single plasmid sequence with isolate-level variation);
\item \textbf{Unique} --- a single dereplication cluster.
\end{itemize}

This partitions the training corpus into 40,871 Major, 24,850 Moderate, and
147,672 Unique PTUs. Within the Moderate category, PTUs contain a median of 2
dereplication clusters (mean 2.33, maximum 18). Within the Unique category,
92.3\% of PTUs comprise a single sequence, while the remaining 7.7\% contain
several near-identical sequences collapsed into one cluster (maximum 84).


\subsection*{Epoch construction as diversity-aware sampling}

Throughout, an epoch denotes one complete pass through the windowed examples
generated from a newly sampled subset of the IMG/PR training set, rather than
through all training sequences. Before windowing, each epoch contained
approximately 166,000 sequence representatives---roughly one quarter of the
training sequences---and the subset and corresponding windowed dataset were
reconstructed at every epoch boundary. For a sequence \(s\), let \(\pi_s\)
denote the marginal probability that \(s\) is included in the sampled subset
before windowing, and define

\[
P(s)=\frac{\pi_s}{\sum_{s'}\pi_{s'}}.
\]

Thus, \(P(s)\) describes the normalized expected sequence composition of an
epoch. It does not imply that sequences within an epoch are sampled
independently, nor does it describe the distribution of windows or nucleotides
presented to the model.

Let \(i\) denote a PTU, \(g(i)\) its category (Major, Moderate or Unique), \(r\) a
possible category value and \(u\) a sub-PTU. The expected sequence composition
factorizes as

\[
P(s)=P(r)\,P(i\mid r)\,P(u\mid i)\,P(s\mid u).
\]

For each Major PTU \(i\), containing \(n_i\) represented sub-PTUs, we defined
the expected allocation

\[
a_i=
\max\!\left[
1,\;
\beta\,\min(n_i,Q_{0.95})^{1/T}
\right],
\]

where \(\beta=0.678\), \(T=1.5\), and \(Q_{0.95}\) is the 95th percentile of
sub-PTU counts among Major PTUs. The resulting allocation was additionally
clipped to the number of sub-PTUs available. The exponent \(1/T=2/3\) increased
allocation with internal diversity but imposed diminishing returns; the floor
ensured representation of every Major PTU, whereas the percentile cap limited
the influence of the most heavily sampled PTUs. For Unique PTUs, let
\(N_{\mathrm{Unique}}\) denote their total number and define

\[
m_{\mathrm{Unique}}
=\left\lfloor f_{\mathrm{Unique}}N_{\mathrm{Unique}}\right\rfloor,
\qquad
q_{\mathrm{Unique}}
=\frac{m_{\mathrm{Unique}}}{N_{\mathrm{Unique}}},
\]

where \(f_{\mathrm{Unique}}=0.564\). The expected PTU-level contribution was
therefore

\[
w_i=
\begin{cases}
a_i, & i\in\mathrm{Major},\\
1, & i\in\mathrm{Moderate},\\
q_{\mathrm{Unique}}, & i\in\mathrm{Unique}.
\end{cases}
\]

Writing \(W_r=\sum_{i:g(i)=r}w_i\) and \(W=\sum_r W_r\), the first two factors
are

\[
P(r)=\frac{W_r}{W},
\qquad
P(i\mid r)=\frac{w_i}{W_r}.
\]

Consequently, Moderate and Unique PTUs are weighted uniformly within their
respective categories, whereas Major PTUs are weighted according to their expected
allocations. The values of \(\beta\) and \(f_{\mathrm{Unique}}\) were chosen to
target an approximately 35\% Major, 15\% Moderate and 50\% Unique composition
among selected representatives; these are target proportions rather than fixed
realized probabilities. Across three independently seeded epoch constructions,
the mean composition was 58,057 Major, 24,848 Moderate and 83,287 Unique
representatives, corresponding to fractions of 0.349, 0.150 and 0.501,
respectively, and a mean total of 166,192 representatives (s.d. 2.9).

Within PTU \(i\), sub-PTUs have equal expected weight,

\[
P(u\mid i)=\frac{1}{n_i},
\]

with \(n_i=1\) for Moderate and Unique PTUs. Representative selection proceeded
by first drawing one sequence uniformly from each dereplication cluster and
then drawing one of the resulting representatives uniformly from each sub-PTU.
Accordingly,

\[
P(s\mid u)
=
\frac{1}{d_u}\frac{1}{|c(s)|},
\]

where \(d_u\) is the number of dereplication clusters in sub-PTU \(u\), and
\(|c(s)|\) is the number of sequences in the dereplication cluster containing
\(s\). This reduces the influence of repeatedly sampled near-identical
sequences without treating every raw sequence as an equally weighted
observation.

Each epoch realized these marginal weights through coupled,
without-replacement sampling. For a Major PTU, write
\(\delta_i=a_i-\lfloor a_i\rfloor\). The realized integer allocation \(K_i\)
was

\[
K_i=
\begin{cases}
\lceil a_i\rceil, & \text{with probability }\delta_i,\\
\lfloor a_i\rfloor, & \text{with probability }1-\delta_i,
\end{cases}
\qquad
\mathbb{E}[K_i]=a_i.
\]

Thus, an allocation of \(a_i=1.2\) yielded two sub-PTU representatives with
probability 0.2 and one with probability 0.8. The required sub-PTUs were then
sampled uniformly without replacement. Each Moderate PTU contributed one
representative. Exactly \(m_{\mathrm{Unique}}\) Unique PTUs were sampled
uniformly without replacement. Thus, Unique-PTU inclusion indicators were not
independent within an epoch, although every Unique PTU had the same marginal
inclusion probability \(q_{\mathrm{Unique}}\). If \(Q_{0.95}=10\), the
expected-allocation ceiling for Major PTUs was 3.15 representatives per epoch,
whereas randomized rounding permitted a maximum realized contribution of four.

At each epoch boundary, representatives within dereplication clusters and
sub-PTUs, randomized Major-PTU allocations, selected Major sub-PTUs, the
Unique-PTU subset, and window phase and window subset were regenerated. These
operations are hierarchical rather than independent. MLM masks were generated
separately on the fly during batch collation.

The resulting sequences were windowed in both orientations, so the
factorization above describes representative selection before the
sequence-length-dependent windowing stage. Every Major and Moderate PTU was
represented in each complete epoch; treating successive epochs as independent
draws, the probability that a given Unique PTU was never selected across the
approximately 102 epochs of training was \((1-q_{\mathrm{Unique}})^{102}\approx
1.7\times10^{-37}\).


\subsection*{Sequence preprocessing, windowing, and tokenization}

\subsubsection*{Sequence sanitization and topology-aware preprocessing}

Each selected sequence was uppercased, and characters outside \{A, C, G, T,
N\} were replaced with N before further processing. Preprocessing then
followed the IMG/PR completeness and topology annotations
(\suppfref{fig:supp-windowing}).

For sequences annotated as putatively complete and as containing a direct
terminal repeat (DTR), the DTR length was identified by a binary search over
candidate terminal overlaps and one terminal copy was removed when detected.
Every sequence annotated as putatively complete was then treated as circular:
1,024 bp from its 3$'$ terminus was prepended and 1,024 bp from its 5$'$
terminus was appended. The extended sequence was randomly offset and trimmed
to a multiple of the 1,024 bp stride before window extraction. This supplied
genuine flanking context to windows crossing the arbitrary linearization
point.

Sequences not annotated as complete were padded with 1,024 placeholder
characters at each end and assigned a uniformly random start offset in
\([0,1{,}024)\). After window extraction, placeholder characters were removed,
leaving terminal windows shorter than the full input length. These
placeholders were used only to position the tiling and never reached the
model; short terminal windows were instead padded with \texttt{[PAD]} during
tokenization.

\subsubsection*{Window generation and strand augmentation}

Processed sequences were tiled into 2,048 bp windows at a stride of 1,024 bp,
giving 1,024 bp overlap between adjacent windows. Each sequence was processed
in both its forward and reverse-complement orientation, and windows from both
orientations were retained. Longer sequences therefore generated
proportionally more training windows. Before window subsampling, an epoch
contained approximately 8.56 million windows.

\subsubsection*{Per-epoch window subsampling}

After shuffling the generated window dataset, a fixed-size 25\% subset was
retained globally without replacement,

\[
m_{\mathrm{window}}=\left\lfloor 0.25\,N_{\mathrm{window}}\right\rfloor,
\]

reducing an epoch to approximately 2.14 million windows. This sampling was
performed over windows, not independently within each sequence. Because the
selected sequences and window offsets were regenerated at each epoch
boundary, the retained windows also varied between epochs. Window subsampling
reduced within-epoch redundancy and bounded the per-process memory required to
materialize each epoch.

\subsubsection*{Tokenization}

Windows were tokenized at single-nucleotide resolution using a seven-token
vocabulary: \texttt{[PAD]} (0), \texttt{[MASK]} (1), \texttt{[UNK]} (2), and
\texttt{a} (3), \texttt{c} (4), \texttt{g} (5), \texttt{t} (6). The tokenizer
lowercased its input, no sequence-level special tokens were added, and
ambiguous N positions mapped to \texttt{[UNK]}. Windows shorter than 2,048
tokens were padded to that length with \texttt{[PAD]}.

Tokenization was performed on the fly during batch collation rather than
materialized in advance; MLM masking was applied separately by the data
collator as described below.


\subsection*{Model architecture, objective, and optimization}

\subsubsection*{Architecture}

\model adapts the GPN convolutional DNA language-model framework
\cite{benegas2023dna} using a ByteNet-style encoder \cite{kalchbrenner2016neural}
of 72 residual blocks with a hidden dimension of 768 and a convolutional
kernel width of 9, totalling 467,904,007 parameters. Each block is a
pre-normalized residual unit applying, in order, layer normalization, GELU, a
linear projection, layer normalization, GELU, a dilated one-dimensional
convolution, layer normalization, GELU, and a second linear projection, with
the block output added to its input. Convolutions are computed at full hidden
width and without bias terms.

The convolutions are \textbf{non-causal}: each position is convolved with
symmetric context on both sides. This differs from the original ByteNet
formulation, in which convolutions are masked to be causal for autoregressive
generation, and is required here because the training objective is masked
prediction rather than next-token prediction.

\subsubsection*{Input representation}

Token identities are encoded as one-hot vectors of width equal to the hidden
dimension, so the first seven channels carry nucleotide identity and the
remainder are zero at input. There is no learned embedding matrix. No
positional encoding is used: position is represented implicitly by the
convolutional stack, which is translation-equivariant by construction --- an
appropriate inductive bias for genomic sequence, where a motif carries the
same meaning wherever it occurs.

\subsubsection*{Dilation and receptive field}

Dilation follows a six-layer cycle of 1, 2, 4, 8, 16, 32, repeated twelve
times across the 72 layers. The resulting receptive field is

\[1 + \sum_{\ell} (k - 1)\,d_\ell \;=\; 1 + 8 \times (1{+}2{+}4{+}8{+}16{+}32) \times 12 \;=\; 6{,}049 \text{ bp},\]

approximately three times the 2,048 bp input window. Every output position
therefore conditions on the entire window rather than on a local
neighbourhood, so the model has full-context access without attention.
Combined with the circularization of complete plasmids, this means that
predictions near the ends of a window are informed by genuine flanking
sequence.

The architecture and masked-prediction workflow are summarized in
\fref{fig:pretraining}g.

\subsubsection*{Pretraining objective}

The model is trained by masked language modelling. \cite{devlin2018bert}
Masking is applied on the fly at batch collation: 15\% of eligible positions
are selected, of which 80\% are replaced by \texttt{[MASK]}, 10\% by a
randomly drawn token, and 10\% left unchanged.

Special tokens are excluded from selection, which has two consequences worth
stating explicitly. \texttt{[PAD]} positions in short terminal windows are
never masked and never supervised. Ambiguous positions, which map to
\texttt{[UNK]}, are likewise never masked and never supervised, so N runs are
neither prediction targets nor a source of spurious loss. The masking rate
therefore applies to nucleotide positions specifically, and the supervised
label set is exactly \{a, c, g, t\}: although the output head is dimensioned
to the full seven-token vocabulary, only the four nucleotide classes are ever
targets.

The prediction head is a two-layer perceptron --- Linear(768 $\rightarrow$
768), GELU, layer normalization, Linear(768 $\rightarrow$ 7) --- and the loss
is cross-entropy over the vocabulary, evaluated only at selected positions;
all other positions carry the ignore index and contribute nothing to the
gradient.

Because the encoder is convolutional and carries no attention mask,
\texttt{[PAD]} tokens in short terminal windows are consumed as ordinary
context by the convolution even though they are excluded from the loss.

\subsubsection*{Optimization}

The model was trained from random initialization with AdamW, a learning rate
of $5 \times 10^{-4}$, and weight decay of 0.01, under a constant schedule
preceded by 2,000 warmup steps. Training used BF16 mixed precision
throughout, distributed data parallelism, and ahead-of-time graph compilation
of the model.

\subsubsection*{Batch schedule and training length}

Per-device batch size was held at 80, and at 76 after the first phase. The
number of devices changed twice over the course of the run: training
proceeded on eight NVIDIA H200 GPUs through step 295,000, then on five, and finally on
four, giving effective batch sizes of 640, 380, and 304 windows in the three
phases. The learning rate was held at $5 \times 10^{-4}$ across all three
phases with no rescaling.

Training completed \textbf{375,000 optimizer steps}. Because the schedule is
constant after warmup, the step count at which training stopped is not a
point on a decay curve --- the released checkpoint is taken from a phase of
constant learning rate, not from a partially completed annealing schedule.

The invariant quantity across the batch-size changes is the amount of data
consumed. Summing the three phases gives approximately 217 million windows,
or about $4.5 \times 10^{11}$ nucleotide tokens. At approximately 2.14
million windows per epoch, this corresponds to roughly 101 resampled epochs,
% Corrected 2026-09-04: 217 / 2.14 = 101.4. The value 102 is used elsewhere as
% an exponent; check supp_tables.tex and the (1-q)^n calculation follow.
each generated from a newly sampled subset of the IMG/PR training set as
defined above.

\subsubsection*{Checkpointing}

Checkpoints were written every 5,000 steps. The released model is the
checkpoint at step 375,000.

\subsubsection*{Pretraining inference}

Inference operates on user-specified genomic positions or pre-defined
windows. For nucleotide prediction at a specified position, a 2,048 bp window
centered on that position is extracted and the central nucleotide is masked.
For embedding extraction, the input table specifies each window and the
requested central region is mean-pooled from the final encoder layer. Each
input is evaluated in its forward and reverse-complement orientation, and the
resulting nucleotide logits or final-layer embeddings are averaged across
orientations. Dense full-sequence maps require an external tiling and
overlap-aggregation step rather than being produced by a single inference
entry point.
